% =====================================================================
\section{Introduction}
\label{sec:intro}
% =====================================================================

\ransac~\cite{Fischler1981} estimates a geometric model from contaminated
correspondences: minimal samples generate candidate models, a score ranks
them, and the best one is kept.  Every widely used score, whether inlier
counting (\ransac), truncated squared error (\msac~\cite{Torr1998}) or the
\gaussunif marginal likelihood (\mlesac~\cite{Torr2000}), depends on a
threshold in residual units, and a user must either set it or estimate the
inlier scale it stands for from contaminated data.

Shekhovtsov~\cite{Shekhovtsov2025} unifies these scores through a single
probabilistic model, the \gaussunif mixture, in which an inlier residual is
Gaussian with scale~$\sigma$ and an outlier is uniform.  \msac is its profile
likelihood, maximized over the latent inlier labels, and \mlesac its marginal
likelihood, summed over them.  Both are parameterized by a decision
threshold $\tau = q\,\sigma$, a fixed quantile of the inlier scale, so tuning
the threshold is tuning~$\sigma$ in other units.  We call these scores
\emph{threshold-based}.  \magsac and \magsacpp~\cite{Barath2019,Barath2020}
were introduced to remove this dependence by ``marginalizing'' the inlier
scale, but their score is numerically equivalent to \mlesac at a fixed
scale~\cite{Shekhovtsov2025}, and their upper bound on~$\sigma$ is a
threshold in disguise.  When the threshold is well tuned, all of these scores
select models equally well; the bottleneck is the threshold itself.

Two lines of work avoid a hand-set threshold.  Adaptive-scale estimators
(\assc~\cite{Wang2004}, \askc~\cite{Wang2010}), SIMFIT~\cite{Heinrich2013}
and \simfitpp~\cite{Edstedt2025} estimate the inlier scale from the
contaminated residuals and then score against it; \simfitpp reduces the
fitting and truncation bias of that estimate with held-out residuals.
\recon~\cite{Raguram2011} seeks consistency in how different models rank the
points.  A-contrario methods (\minpran~\cite{Stewart1995},
\orsa/\acransac~\cite{Moisan2004,Moisan2012}) need no inlier scale: they rank
inlier sets by a significance test against a uniform background model, whose
domain they take from the image.  None of them fully parameterizes the
\gaussunif mixture that the threshold-based scores share.

The profile-marginal score of Pritts et~al.~\cite{Pritts2026} does, for the
inlier half.  It reverses the usual order of inference: instead of
eliminating the labels at a fixed scale, it integrates out~$\sigma$ under a
conjugate \invgamma prior for each fixed partition into inliers and
outliers, and then maximizes over partitions.  The integral has a closed
form, one \sortsweep finds each candidate's best partition, and the inlier
scale is absent from the score.  The outlier half is still fixed by hand:
the half-width~$a$ of the uniform outlier distribution is a length in
residual units.  The score is far less sensitive to~$a$ than threshold-based
scores are to their threshold, because its implicit inlier threshold depends
on~$a$ only logarithmically~\cite{Pritts2026}, but $a$ must still be chosen
for each residual type, coordinate normalization and image size, it is
cross-validated like a threshold, and it is why the score's ranking of
models still depends on the units of the residuals.

We parameterize the outlier half of the mixture from the data as well.  The
key observation is that profiling over partitions gives every candidate
model, at no cost, a natural complement: the correspondences its best
partition labels as outliers.  The upper tail of their residuals is exactly
what~$a$ describes, the extent of the outlier support, and extreme value
theory gives that tail a universal form.  Above a high threshold, the
exceedances of almost any distribution follow a generalized Pareto
distribution~\cite{Pickands1975,Balkema1974}, whose endpoint is the support
boundary.  We fit it to the complement by peaks over
threshold~\cite{Coles2001,Hosking1987}, with the threshold set as a quantile
of the complement, and replace the half-width by the estimated endpoint.
Both components of the mixture are then fitted to the partition they
describe: the inlier scale is marginalized and the outlier support is
estimated.  With the non-informative prior, the score contains no parameter
with units of length.  It is, to our knowledge, the first thresholdless
\ransac score: no threshold, no inlier scale and no outlier width is set,
and its ranking of models is provably invariant to the units of the
residuals.  Its one remaining choice, the threshold quantile, is
dimensionless, and the threshold stability of the generalized Pareto tail
makes the score insensitive to it.

\paragraph{Relation to \evsac.}
The closest work is \evsac by Fragoso et~al.~\cite{Fragoso2013evsac}, which
also brings extreme value theory into \ransac.  The two methods differ in
what they model, at which stage, and to what end.
\evsac models the \emph{matching scores} of a nearest-neighbour feature
matcher, the descriptor distances of each correspondence, as a mixture of a
distribution for correct matches and an extreme value distribution for
incorrect ones, fitted once per image pair and before any model exists.  It
turns this mixture into a confidence per correspondence that guides
\emph{sampling}, so that good hypotheses are generated sooner; the score
that ranks the hypotheses, and its inlier threshold, are left unchanged.
We model the \emph{geometric residuals} of each candidate model, not
descriptor distances, and we change the \emph{score}, not the sampler.  Our
extreme value fit is recomputed for every candidate and every partition,
from the complement that profiling supplies, and it removes the last
residual-unit parameter of the score.  It needs no matching scores, so it
applies unchanged to learned matchers and to any problem with residuals.
The choice of extremes also differs, and it matters.  \evsac models minimum
distances, the natural extremes of nearest-neighbour search.  For the
outlier residuals, the minima are the wrong extremes: they lie where the
inliers are, an estimate built from them adapts to the inliers in the
complement, and the score collapses to its smallest partition
(\cref{sec:why-maxima}).  The upper tail is formed by the gross outliers far
from the model, which do not move as the partition grows.  The two methods
are complementary, and \evsac's sampler can feed our score.
\TODO{check the description of EVSAC against the paper (Gamma/GEV mixture,
fitted from the 2nd..k-th neighbours; confidences used for guided sampling)}

\paragraph{Contributions.}
\begin{enumerate}
	\item We show that the profile-marginal score supplies, for every
	      candidate, a complement of putative outliers whose upper tail
	      identifies the outlier support, and we estimate that support by
	      peaks over threshold.  This fully parameterizes the \gaussunif
	      mixture from the data and gives a thresholdless \ransac score.
	\item We prove that its ranking of models is invariant to the units of
	      the residuals, and that its one remaining choice, the threshold
	      quantile, does not change the estimated support.
	\item We show why the upper tail, not the minima, is the right extreme:
	      a minima-based estimate adapts to the inliers in the complement and
	      collapses the partition.
	\item The fitted tail shape is a free diagnostic of the uniform outlier
	      model.  \TODO{experimental contribution on real pools}
\end{enumerate}
